\scanpage{22}
\deleted{\unsure{car} $U_{\varphi_K}$ est dans $\underline b$, \unsure{ceci implique}}
Soit $f$ fonction borélienne bornée sur $M$ nulle à l'$\infty$, \deleted{l'opérateur $U_f$ est dans} \deleted{Si $f$ est à support compact, $|U_f|$ \inserted{$U_{|f|}$} est majoré par $U_g$ avec $g\in C_0(\Omega)$.}

\unsure{En effet}. $|f|$ est majorée par une $g\in C_0(M)$, d'où $0\leq|U_f|\leq U_g\in\underline b$, d'où $U_f\in\underline b$, i.e. $U_f$ est de la forme $U_h$, $h\in C_0(M)$. En particulier, cela est vrai si $f$ est de la \deleted{forme $\varphi_K$, $K$ compact $\subset\Omega$.} forme $\varphi_K$, $K$ compact $\subset M$. Donc il existe alors une partie compacte et ouverte $K'$ telle que $U_{\varphi_K}=U_{\varphi_{K'}}$, on a $K\subset K'$.

Prouvons que $K$ \unsure{sont ouvertes et fermées|est ouverte et fermée}, \inserted{soit $E=U_{\varphi_K}$.} Tout opérateur \inserted{$A$} de $\mathcal A$ qui permute à $\underline c$, et tel que $EAE=A$, est dans $\underline c$, et \unsure{est} dans $C(K)$. Donc $\unsure{K}$ est le spectre de l'algèbre formée des opérateurs précédents, algèbre qui est \deleted{\unsure{commutative}} \unsure{fblt} fermée. Donc $M$ est Stonien.

\uline{Corollaire 1} Dans le cas où $\unsure{\underline a|\underline c}$ est l'idéal de définition d'une trace \inserted{\unsure{s. finie|semi-finie}} $T$, $\underline c$ une sous-algèbre commutative \unsure{uniformt} fermée de $\underline b$ = adh. de $\underline a$, $\underline c=C_0(M)$, alors on peut définir sur $\Omega$ une mesure positive $\unsure{\mu|\rho}$ telle que si $f\in C_0(\Omega)$, on ait $f\in\underline a$ si et seulement si $f$ est sommable pour $\unsure{\mu|\rho}$. Alors $T(f)=\int f\,d\unsure{\mu|\rho}$. Si $T$ est \inserted{fidèle, $\mathop{\rm supp}\unsure{\mu|\rho}=\Omega$}
\mspage{23}
